\documentclass[a4paper]{article}

\usepackage[spanish]{babel}
\usepackage{graphicx}
\usepackage{amsmath, amssymb}
\usepackage[margin=2cm]{geometry}
\usepackage{fancyhdr}
\usepackage{enumerate}
\usepackage[shortlabels]{enumitem}
\usepackage{parskip}
\usepackage[most]{tcolorbox}
\usepackage[hidelinks]{hyperref}
\usepackage{float}


% cabecera
\pagestyle{fancy}
\fancyhead[l]{Daniel Soto}
\fancyhead[c]{Comando útiles Git-GitHub}
\fancyhead[r]{\today}
\fancyfoot[c]{\thepage}
\renewcommand{\headrulewidth}{0.2pt} % linea horizontal


\begin{document}

\textbf{Cómo clonar un repositorio}

\begin{verbatim}
git clone git@github.com:Daniel-534/AprenderGit.git
\end{verbatim}

\textbf{Borrar Caché total:} Este caso surge del siguiente problema:

Es posible que en cierto momento, aparezcan archivos que no necesitamos en nuestro repositorio, entonces es lo que incluímos en el \texttt{.gitignore}. Pero si ya tenemos los archivos en el GitHub y después actualizamos el \texttt{.gitignore}, notaremos que no se van a desaparecer esos archivos no deseados. Aquí lo que se debe hacer es una eliminación total del caché de manera recursiva, para que pueda entrar a todas las carpetas. El comando que se debe ejecutar es:

\begin{verbatim}
git rm -r --cached .
\end{verbatim}


\end{document}
